\section{Experimental Procedures and Data Analysis}

\subsection{Cold-Atom Preparation: MOT and Optical Molasses}

Each experimental cycle begins with loading atoms into the magneto-optical trap (MOT), using the operating conditions previously identified as optimal\CCgen{capitolo 2}{in tutto il capitolo non c'e' una sola tabella}{Contato: 21 figure, 104 equazioni, \textbf{zero tabelle}. Tutti i parametri operativi vivono nella prosa o in sette elenchi puntati (73 voci in tutto). Il risultato e' che per sapere con quali impostazioni e' stata presa una misura bisogna ricostruirla leggendo pagine diverse, e infatti proprio da qui nasce l'ambiguita' sull'ampiezza AOM, 475 in una sezione e 480 in un'altra, che nessuno dei due punti chiarisce. Una tabella sola, con la configurazione finale (detuning, ampiezze e potenze AOM, durata delle rampe, correnti di compensazione e campi corrispondenti, waist, potenze ai vari stadi), risolve insieme il problema di lettura e quello di coerenza, e diventa il riferimento che il capitolo 3 puo' citare invece di ripetere i numeri.} for this apparatus. The MOT is run with an axial magnetic-field gradient of $B' = 9.32\,\mathrm{G/cm}$, corresponding to a current of $8\,\mathrm{A}$ in the quadrupole coils, a cooling-light detuning of $\Delta = -13\,\mathrm{MHz}$, and all compensation-coil currents set to zero, the loading stage lasts $5\,\mathrm{s}$.

At the end of the MOT phase, the quadrupole magnetic field is switched off and the cloud enters the optical molasses stage.\CCthesis{le citazioni calano man mano che il contenuto diventa tuo}{Contate per capitolo: \textbf{71} citazioni nel capitolo 1, \textbf{18} nel 2, \textbf{5} nel 3. La curva scende proprio dove il confronto con la letteratura conta di piu': nel capitolo 1 si riassumono cose note e si cita molto, nel capitolo 3 si presentano numeri nuovi e non si cita quasi nulla. Andrebbe quasi rovesciata. Non serve gonfiare il capitolo 1: serve che ogni risultato del capitolo 3 e ogni tecnica del capitolo 2 abbiano accanto il riferimento con cui si confrontano, cosa che il confronto con Gerz nel capitolo 3 mostra di saper fare bene. Questa nota guarda l'andamento sull'intera tesi; le due note di capitolo sulla stessa cosa, una nel capitolo 2 e una nel capitolo 3, entrano nel dettaglio di ciascuno.} In the absence of the trapping field, the atoms are no longer confined and the cloud starts to expand. At the same time, the atomic temperature decreases further, while the cloud falls under gravity and remains subject to a velocity-dependent damping force generated by the cooling light.

The temperature achieved during the molasses phase scales as $I/|\Delta|$, where $I$ is the laser intensity and $\Delta$ is the detuning(Sec.~\ref{sec:ODT}). However, this relation remains valid only as long as $|\Delta|$ is small enough to avoid a significant reduction in the optical scattering rate, $\Gamma_p \propto 1/\Delta^2$. Furthermore, the detuning must not be so large as to induce a significant excitation probability for the neighboring $F=2 \rightarrow F'=2$ transition.

A second requirement for efficient molasses cooling is the cancellation of residual magnetic fields. For this reason, $5\,\mathrm{ms}$ after the MOT field is switched off, the compensation-coil currents are adjusted in order to compensate for the Earth's magnetic field and for other stray fields present in the laboratory. The cooling-light detuning is then changed from the MOT value to the molasses value through a 20-step ramp with total duration $5\,\mathrm{ms}$. This ramp starts $50\,\mu\mathrm{s}$ after the compensation currents are set to their molasses values.

\subsection{Atom Imaging}

The atomic cloud is observed inside the vacuum chamber by means of two monochrome CMOS cameras:

\begin{itemize}
    \item a Teledyne FLIR BFS-U3-04S2M-CS camera, based on the Sony IMX287 sensor, with a resolution of $720\times540$ pixels and a pixel size of $6.9\,\mu\mathrm{m}\times6.9\,\mu\mathrm{m}$;
    \item an Allied Vision Manta G-419 NIR camera, based on the CMOSIS/ams CMV4000 sensor, with a resolution of $2048\times2048$ pixels and a pixel size of $5.5\,\mu\mathrm{m}\times5.5\,\mu\mathrm{m}$.
\end{itemize}

The FLIR camera is used mainly during the optimization of the molasses stage. The Manta camera, owing to its higher sensitivity at $780\,\mathrm{nm}$ and its smaller pixel pitch, is used for ODT loading measurements, where resolving the small trapped cloud is more demanding.

Both cameras are equipped with a Thorlabs MVL50M23 adjustable objective, with focal length $50\,\mathrm{mm}$, minimum working distance $200\,\mathrm{mm}$, and minimum magnification of about $0.33$. The imaging system is focused so that the HCPCF inside the chamber, and therefore also the atomic cloud located at nearly the same distance, is sharply imaged on the sensor. In both cases the acquisition is triggered externally by the FPGA control system.

\subsection{FLIR Camera Calibration}

A calibration of the imaging system is required in order to convert distances measured on the sensor into physical lengths in the object plane. This is needed for the temperature measurements, where the cloud width extracted from the images must be expressed in absolute units. Starting from the physical pixel pitch of the FLIR sensor, $6.9\,\mu\mathrm{m}$, the optical magnification $M$ and the corresponding object-plane pixel size $C$ were estimated through three independent methods.

\begin{enumerate}

\item \textbf{Direct imaging of a known microstructure}

The first method uses the hollow-core photonic crystal fiber (HCPCF) as a reference object, exploiting its known physical diameter of $316\,\mu\mathrm{m}$. The fiber is imaged directly to determine its apparent diameter on the sensor, which is then compared with its actual physical diameter. To extract the apparent diameter, the image intensity is summed along the columns, and the resulting one-dimensional profile is fitted with a rectangular function.
\CCn{The fitted image corresponds to an apparent diameter of $8 \pm 2$ pixels.}{Vale la pena dire che questa prima calibrazione, da sola, non e' utilizzabile: un diametro di $8\pm2$ pixel porta un'incertezza del 25 per cento, che si trasferisce tale e quale su $C$ ($39.5\pm9.9~\mu$m/px) e su $M$ ($0.17\pm0.04$). E' venti volte peggio del metodo del righello, e infatti quello e' il metodo che adotti. Serve come controllo d'ordine di grandezza, non come misura, e presentarla per prima e con lo stesso peso delle altre due fa pensare il contrario. Basta una frase che dica che il metodo 1 e' un controllo grossolano, limitato dal fatto che l'oggetto di riferimento copre appena otto pixel.} The resulting object-plane calibration is therefore
\begin{equation}
    C = \frac{316\,\mu\mathrm{m}}{(8 \pm 2)\,\mathrm{px}} = 39.5 \pm 9.9\,\mu\mathrm{m/px},
\end{equation}
which gives a magnification
\begin{equation}
    M = \frac{6.9}{39.5 \pm 9.9} = 0.17 \pm 0.04.
\end{equation}

\item \textbf{Optical calibration by reflected ruler imaging}

The second method relies on imaging a macroscopic reference target placed at the same optical distance as the atomic cloud. A ruler is observed through a mirror after the camera has already been focused on the atoms. The ruler is translated until its reflected image becomes sharp, ensuring that the target lies on the same focal plane as the cloud.

In order to reduce the uncertainty and average over possible local imperfections, multiple measurements are performed. Specifically, five independent measurements are taken by repositioning the ruler at focus each time, finding that a physical interval of $2\,\mathrm{cm}$ corresponds to $471 \pm 2\,\mathrm{px}$. This procedure yields an effective pixel size
\begin{equation}
    C = 43.4 \pm 0.3\,\mu\mathrm{m/px},
\end{equation}
corresponding to a magnification
\begin{equation}
    M = \frac{6.9}{43.4 \pm 0.3} = 0.159 \pm 0.001.
\end{equation}
\item \textbf{Calibration from MOT displacement in a bias magnetic field}

As a further cross-check, the imaging scale can be estimated from the displacement of the MOT induced by a controlled bias magnetic field. When a uniform field is applied along the vertical direction by means of the $y$-axis compensation coils, the zero of the quadrupole field is shifted, and the cloud follows the new equilibrium position.

Experimentally, the cloud displacement varies linearly with the current applied to the $y$-axis compensation coils. Furthermore, the beam trajectory on the sensor shows a linear correlation between the horizontal and vertical displacements. Converting the fitted values from the plot into px/A, the measured slopes are
\begin{equation}
    x_{\mathrm{shift}} = 20.3 \pm 0.5\,\mathrm{px/A},
\end{equation}
\begin{equation}
    y_{\mathrm{shift}} = 58.6 \pm 0.6\,\mathrm{px/A}.
\end{equation}

To convert this result into a physical displacement, one needs the bias field generated per unit compensation current and the restoring gradient of the quadrupole field. Modeling the compensation coils as a cubic structure of side $L = 645\,\mathrm{mm}$ with $N = 100$ windings gives
\begin{equation}
    b = 1.43\,\mathrm{G/A}.
\end{equation}
The quadrupole gradient is estimated from the coil geometry. For a current of $8.4\,\mathrm{A}$, 128 windings, effective radius $R_{\mathrm{eff}} = 62\,\mathrm{mm}$, and effective separation $a_{\mathrm{eff}} = 162\,\mathrm{mm}$, the vertical gradient is
\begin{equation}
    G_{y,\mathrm{tot}} = 0.571\,\mathrm{G/mm}.
\end{equation}

The equilibrium displacement per unit current is given by the ratio between the applied bias field and the magnetic-field gradient. Dividing this physical shift by the measured vertical displacement in pixels gives the calibration factor:
\begin{equation}
    C = \frac{b}{G_{y,\mathrm{tot}}\,y_{\mathrm{shift}}}
    = \frac{1.43\,\mathrm{G/A}}{0.571\,\mathrm{G/mm}\times (58.6 \pm 0.6)\,\mathrm{px/A}}
    = 42.7 \pm 0.4\,\mu\mathrm{m/px}.
\end{equation}

If, instead, the empirically measured values previously obtained for this setup are used~\cite{Marchesini} namely $b = 1.5\,\mathrm{G/A}$ and $G = 5.6\,\mathrm{G/cm}$ at $8.4\,\mathrm{A}$, one obtains
\begin{equation}
    C = \frac{1.5\,\mathrm{G/A}}{5.6\,\mathrm{G/cm}\times (58.6 \pm 0.6)\,\mathrm{px/A}} = 45.7 \pm 0.5\,\mu\mathrm{m/px}.
\end{equation}

\CCerr{Both estimates are consistent with the result obtained from the ruler calibration.}{I numeri dicono il contrario per la seconda. Righello $43.4\pm0.3$; magnetica geometrica $42.7\pm0.4$, che dista $1.4\sigma$ ed e' quindi compatibile; magnetica con i valori di Marchesini $45.7\pm0.5$, che dista $\mathbf{3.9\sigma}$. E le due stime magnetiche fra loro distano $\mathbf{4.7\sigma}$. La causa e' visibile nelle formule: l'incertezza citata propaga solo $y_{\mathrm{shift}}$, mentre $b$ e $G$ entrano al numeratore e al denominatore e non hanno barra d'errore, pur essendo l'una stimata dalla geometria delle bobine e l'altra presa da una misura precedente. Sono loro a dominare. Due strade: propagare anche le incertezze di $b$ e $G$, e allora molto probabilmente tutto torna compatibile e l'affermazione diventa vera; oppure, se restano cosi', dire che la terza stima e' in disaccordo e discuterne il perche'. Cosi' com'e' il testo afferma una compatibilita' che i suoi stessi numeri smentiscono.} In the following analysis, the value provided by the ruler method is adopted.

A summary of the three calibration methods is shown in Fig.~\ref{fig:calibration_summary}.

\begin{figure}[h!]
    \centering
    % Prima colonna a sinistra: Figura (a) in verticale
    \begin{minipage}[c]{0.15\textwidth}
        \centering
        \includegraphics[width=\linewidth]{Cap_2_figures/calibration_fiber_photo.png}\\[4pt]
        \textbf{(a)}
    \end{minipage}
    \hfill
    % Seconda colonna centrale: Figura (b) in alto e Figura (c) sotto di essa
    \begin{minipage}[c]{0.20\textwidth}
        \centering
        \includegraphics[width=\linewidth]{Cap_2_figures/calibration_fiber.jpeg}\\[4pt]
        \textbf{(b)}
        
        \vspace{12pt} % Spazio verticale tra (b) e (c)
        
        \includegraphics[width=\linewidth]{Cap_2_figures/calibration_ruler.png}\\[4pt]
        \textbf{(c)}
    \end{minipage}
    \hfill
    % Terza colonna a destra: Figura (d) molto più grande e orizzontale
    \begin{minipage}[c]{0.60\textwidth}
        \centering
        \includegraphics[width=\linewidth]{Cap_2_figures/calibration_magnetic.png}\\[4pt]
        \textbf{(d)}
    \end{minipage}
    
    \vspace{10pt} % Spazio tra le immagini e la didascalia
    \caption{Summary of the spatial calibration methods used for the imaging system. \textbf{(a)} Side-view image of the HCPCF used as the reference object in Method 1. \textbf{(b)} Cross-sectional view of the same fiber, with the nominal diameter of $316\,\mu\mathrm{m}$. \textbf{(c)} Ruler used for the reflected-image calibration of Method 2. \textbf{(d)} MOT position as a function of the current applied to the $y$-axis compensation coils in Method 3; the linear fit yields a slope of about $58\,\mathrm{px/A}$.}
    \label{fig:calibration_summary}
\end{figure}

\end{enumerate}

\subsection{Temperature Measurements}

\CCerr{Once the imaging calibration is known, the cloud size extracted from the images can be expressed in physical units and used to determine the temperature from ballistic expansion.}{Vale la pena scrivere qui che la calibrazione entra al QUADRATO, perche' e' il punto in cui il disaccordo fra i tre metodi segnalato poco sopra smette di essere una questione interna e diventa un errore sulla temperatura. La legge di fit da' $\sigma_x^2$ in pixel al quadrato, quindi $T\propto C^2$: uno scarto del $5.3$ per cento fra il righello ($43.4~\mu$m/px) e la stima magnetica con i valori di Marchesini ($45.7$) diventa il $10.9$ per cento su $T$. Sulla temperatura della melassa, $T=(8.42\pm0.17)~\mu$K, sono $0.92~\mu$K, cioe' \textbf{$5.4$ volte l'incertezza citata}. In altre parole il bilancio degli errori sulla temperatura e' dominato da un sistematico che non compare nel bilancio: la barra riportata e' quella del fit e basta. Due cose da fare, entrambe brevi: dire esplicitamente che $T\propto C^2$, e propagare l'incertezza della calibrazione adottata dentro l'incertezza di $T$. Con il solo righello, $43.4\pm0.3$, il contributo e' l'$1.4$ per cento e la barra cresce poco; e' se si tiene aperto il disaccordo fra metodi che diventa il termine dominante.} These measurements are used to optimize the molasses parameters by identifying the conditions that minimize the final temperature.

The method is based on the free expansion of the cloud after release. For a non-interacting gas, the position of an atom along one direction evolves as
\begin{equation}
    x(t) = x_0 + v_x t.
\end{equation}
Assuming that the initial position and velocity are statistically independent, the variance at time $t$ is
\begin{equation}
    \sigma_x^2(t) = \sigma_{x,0}^2 + \sigma_{v_x}^2 t^2,
\end{equation}
where $\sigma_{x,0}$ is the initial spatial width. In the center-of-mass frame, the equipartition theorem gives
\begin{equation}
    \sigma_{v_x}^2 = \langle v_x^2 \rangle = \frac{k_B T_x}{m}.
\end{equation}
The fitting law then becomes
\begin{equation}
    \sigma_x^2(t) = \sigma_{x,0}^2 + \frac{k_B T_x}{m} t^2,
    \label{eq:ballistic_expansion}
\end{equation}
with $m$ the mass of $^{87}\mathrm{Rb}$, $m \approx 1.443\times10^{-25}\,\mathrm{kg}$.

The experimental procedure is as follows:

\begin{enumerate}
    \item A MOT is loaded and followed by the molasses stage. The molasses parameters and duration are the quantities under optimization.
    \item The cooling light is switched off, allowing the cloud to expand freely and fall under gravity. At each selected time of flight, the FLIR camera acquires a fluorescence image of the cloud.
    \item A background frame is then acquired after a sufficiently long delay, so that the cloud has completely left the observation region.
    \item For each time of flight, the corrected image is obtained as
    \begin{equation}
        I(t) = I_{\mathrm{raw}}(t) - I_{\mathrm{bg}}(t).
    \end{equation}
    \item Each corrected image is fitted with a two-dimensional Gaussian, yielding the widths $\sigma_x$ and $\sigma_y$.
    \item Equation~\ref{eq:ballistic_expansion} is fitted independently along the two axes to extract $T_x$ and $T_y$. The temperature quoted in the following is the average of the two values.
\end{enumerate}
The imaging process is destructive because the light flash heats the sample, resulting in its loss. Therefore, each time-of-flight measurement cannot be done sequentially, but requires a new experimental run from start to finish.

\subsection{ODT Loading Measurements}
\label{subsec:odt_loading_measurements}
The loading performance of the optical dipole trap is characterized using the Manta camera. Two quantities are extracted from these measurements: the loading efficiency $\eta$ and the integrated fluorescence signal associated with the trapped cloud, which is proportional to the number of captured atoms.

Because only a small fraction of the molasses atoms, of the order of $10^{-3}$, is transferred into the ODT, the trapped-atom signal is much weaker than the fluorescence produced by the untrapped cloud. For this reason, the analysis is based on differential imaging. 

\begin{figure}[h!]
    \centering
    \includegraphics[width=0.8\linewidth]{Cap_2_figures/loading_free_fall.png}
    \caption{Schematic representation of the experimental sequence for loading and guiding atoms. The process consists of four stages: (1) preparation of the cold atomic cloud, (2) loading the atoms into the far-off-resonance trap (FORT) using a trapping beam, (3) free-falling of the atomic cloud under gravity, and (4) guiding of the optically trapped atoms into the hollow-core fiber (HCF), while the untrapped atoms fall and disperse outside the trap. Credits~\cite{Wang2025}.}
    \label{fig:loading_free_fall}
\end{figure}

The measurement sequence is the following (illustrated in Fig. \ref{fig:loading_free_fall}):

\begin{enumerate}
    \item Two otherwise identical experimental cycles are performed in succession, one with the ODT on from the MOT loading stage onward and one with the ODT off.
    \item \CCn{After a time of flight of $30\,\mathrm{ms}$, long enough for the untrapped atoms to expand away from the ODT region, fluorescence images $I_{\mathrm{on}}$ and $I_{\mathrm{off}}$ are acquired with the Manta camera.}{La scelta e' giusta e il numero che la giustifica si ricava dai tuoi stessi dati, quindi conviene metterlo invece di lasciare ``long enough''. Con $T=8.4~\mu$K la velocita' quadratica media e' $\sigma_v=\sqrt{k_BT/m}=28$~mm/s, quindi in $30$~ms la nuvola libera si allarga di $\sigma_v t=0.85$~mm, mentre la caduta e' $gt^2/2=4.4$~mm: il centro di massa si sposta di circa $5.2$ larghezze. E' questo il margine che rende separabili le due componenti, ed e' anche il criterio con cui si sceglierebbe un tempo diverso se la temperatura cambiasse, visto che la separazione in unita' di larghezza vale $gt/2\sigma_v$ e cresce quindi linearmente nel tempo (la caduta va come $t^2$, la larghezza come $t$): a $10$~ms sarebbe $1.7$, che non basterebbe. Una mezza riga con $4.4$~mm contro $0.85$~mm dice al lettore perche' $30$~ms e non $10$.}
    \item Each on/off pair is repeated five times.
    \item The images are cropped to a pixel region of interest centered on the cloud. This excludes both the visible HCPCF structure and the peripheral parts of the image containing only background.
    \item For each condition, the cropped images are averaged:
    \begin{equation}
        \bar I_{\mathrm{on}} = \langle I_{\mathrm{on}} \rangle,
        \qquad
        \bar I_{\mathrm{off}} = \langle I_{\mathrm{off}} \rangle.
    \end{equation}
    \item The averaged images are integrated along one direction to obtain one-dimensional profiles with improved signal-to-noise ratio.
    \item The total number of atoms in the molasses drifts slowly over time because of variations in vacuum pressure and thermal changes in the optical components. As a consequence, the on and off profiles cannot be subtracted directly: even a small mismatch in the broad molasses background would produce a residual comparable to the ODT signal itself.
    
    To correct for this effect, the off profile is rescaled by a factor $k$ chosen so that the difference between the two profiles vanishes outside the ODT region. \CCerr{An exclusion window centered at the ODT position, $x_{\mathrm{ODT}} \pm R_{\mathrm{window}}$, is defined.}{$R_{\mathrm{window}}$ non riceve mai un valore, e non e' un dettaglio: e' il parametro che decide $k$, e $k$ entra due volte nell'efficienza di caricamento, al numeratore attraverso la sottrazione e al denominatore in modo esplicito. Conviene anche dire da che parte sbaglia, perche' il verso e' univoco. Se la finestra e' troppo stretta, parte del segnale dell'ODT finisce nella regione ``outside'': $A_{\mathrm{on}}^{\mathrm{out}}$ si gonfia, $k$ risulta troppo grande, la sottrazione $\bar I_{\mathrm{on}}-k\bar I_{\mathrm{off}}$ toglie troppo e $A_{\mathrm{ODT}}^{\mathrm{G}}$ si riduce; nello stesso tempo il denominatore $k\,A_{\mathrm{off}}^{\mathrm{G}}$ cresce. Numeratore giu' e denominatore su: \textbf{una finestra troppo stretta sottostima $\eta$}, sempre. Se e' troppo larga, $k$ viene stimato sulle code e il rumore su $\eta$ cresce senza distorsione. Servono quindi il valore usato e una riga che mostri che il risultato non dipende dalla scelta, per esempio $\eta$ per due o tre valori di $R_{\mathrm{window}}$: e' la stessa verifica di robustezza che si chiede a una soglia, e qui costa una tabella di tre righe.} Outside this interval, a linear baseline $L_i(x)$ is estimated for each profile and the residual integrated signal is computed as
    \begin{equation}
         A_i^{\mathrm{out}} = \int_{\mathrm{outside}} \left[\bar{I}_i(x) - L_i(x)\right]\,dx,
         \qquad i \in \{\mathrm{on, off}\}.
    \end{equation}
    The normalization factor is then
    \begin{equation}
        k = \frac{A_{\mathrm{on}}^{\mathrm{out}}}{A_{\mathrm{off}}^{\mathrm{out}}}.
    \end{equation}
    
    \item The differential profile attributed to the trapped atoms is therefore
    \begin{equation}
        \Delta \bar I(x) = \bar I_{\mathrm{on}}(x) - k\,\bar I_{\mathrm{off}}(x).
    \end{equation}
    
    \item A Gaussian plus linear background is fitted both to the off profile and to the differential profile. If $A_{\mathrm{off}}^{\mathrm{G}}$ and $A_{\mathrm{ODT}}^{\mathrm{G}}$ denote the Gaussian areas extracted from the two fits, the loading efficiency is defined as
    \begin{equation}
        \eta = \frac{A_{\mathrm{ODT}}^{\mathrm{G}}}{k\,A_{\mathrm{off}}^{\mathrm{G}}}.
    \end{equation}
\end{enumerate}


\subsection{Measurement of the Number of Atoms}
\label{subsec:meas_n_atoms}
The loading efficiency can be obtained from relative fluorescence signals and therefore does not require an absolute calibration of the camera response. Estimating the absolute number of trapped atoms requires converting the recorded counts into a known number of detected photons.

To this end, the camera is calibrated against a beam of known optical power coming from the MOT light. The beam is directed onto the sensor after passing through a series of attenuators and through the vacuum-chamber viewports. 
The total transmission of the optical path is written as the product of the individual transmission factors,
\begin{equation}
    T_{\mathrm{tot}} = T_1\,T_{\mathrm{vac}}\,T_2\,T_3 \approx 0.01272.
\end{equation}
During the calibration, the AOM amplitude is set to 50, corresponding to an input power of $22\,\mu\mathrm{W}$ after background subtraction. The optical power reaching the camera is therefore
\begin{equation}
    P_{\mathrm{cam}} = 22\,\mu\mathrm{W}\times T_{\mathrm{tot}} \approx 279.8\,\mathrm{nW}.
\end{equation}

\begin{wrapfigure}{r}{0.5\textwidth}
    \centering
    \includegraphics[width=\linewidth]{Cap_2_figures/manta_calib.png}
    \caption{Linear fit of the total camera counts as a function of the exposure time $t_{\mathrm{exp}}$. The slope $m$ extracted from the fit is used to calculate the photon conversion factor $K$.}
    \label{fig:calibration_manta}
\end{wrapfigure}

    The camera response is determined by varying the exposure time~$t_{\mathrm{exp}}$ while keeping~$P_{\mathrm{cam}}$ fixed. For each image, the beam profile is fitted in two dimensions and the integrated counts are extracted. Denoting by~$C$ the total number of recorded counts, the response is assumed to be linear:
\begin{equation}
    C = m\,t_{\mathrm{exp}} + q,
\end{equation}
where $m$ is the count rate. Since the incident photon flux is given by $P_{\mathrm{cam}}/(hc/\lambda)$, the conversion factor from counts to detected photons is
\begin{equation}
    K = \frac{P_{\mathrm{cam}}\lambda}{m h c},
\end{equation}
where $h$ is Planck's constant, $c$ the speed of light, and $\lambda$ the laser wavelength. The resulting linear fit is shown in Fig.~\ref{fig:calibration_manta}, from which the conversion factor $K$ is extracted.

For an ODT fluorescence image, the integrated residual signal $C_{\mathrm{tot}}$ can then be converted into the number of detected photons:
\begin{equation}
    N_{\mathrm{det}} = K\,C_{\mathrm{tot}}.
\end{equation}

In order to infer the number of trapped atoms, the fluorescence contribution of a single atom during the imaging pulse must be estimated. The scattering rate is described by the standard two-level expression
\begin{equation}
    R_{\mathrm{sc}} =
    \frac{\Gamma}{2}
    \frac{I/I_{\mathrm{sat}}}{1 + I/I_{\mathrm{sat}} + 4(\Delta/\Gamma)^2},
\end{equation}
where $\Gamma$ is the natural linewidth, $I$ is the total imaging intensity, $I_{\mathrm{sat}}$ is the saturation intensity, and $\Delta$ is the laser detuning. Considering a single MOT beam with a power of $P = 14\,\mathrm{mW}$ and a waist of $W = 12.5\,\mathrm{mm}$, the peak intensity is $I_0 = 2P/(\pi W^2) \approx 5.70\,\mathrm{mW/cm^2}$. Using the ideal two-level saturation intensity for the $^{87}$Rb cooling transition ($I_{\mathrm{sat}} = 1.67\,\mathrm{mW/cm^2}$), \CCn{the saturation ratio is evaluated as $s_0 = I_0/I_{\mathrm{sat}} \approx 3.4$}{Qui c'e' un salto fra la formula e il numero. La formula di $R_{\mathrm{sc}}$ chiede l'intensita' TOTALE che vede l'atomo, come scrivi due righe sopra (``$I$ is the total imaging intensity''), ma il valore inserito e' il picco di UN SOLO fascio del MOT. Se l'immagine viene presa con i sei fasci accesi, l'intensita' e' diverse volte maggiore e $s_0$ passa da $3.4$ a un valore dell'ordine di 20. Non e' drammatico, perche' $R_{\mathrm{sc}}$ satura: a risonanza si passa da $s/(1+s)=0.77$ a $0.95$, cioe' circa il 24 per cento in piu' di fotoni per atomo e altrettanti atomi in meno. Ma va deciso e detto: o si dichiara che l'imaging usa un fascio solo, o si mette l'intensita' totale. Verificato anche il resto del conto e torna: $2\times14/(\pi\cdot1.25^2)=5.70~\mathrm{mW/cm^2}$ e $5.70/1.67=3.41$.}.

For an imaging pulse of duration $t_{\mathrm{img}}$, each atom scatters on average $R_{\mathrm{sc}} t_{\mathrm{img}}$ photons. Only a fraction of this fluorescence reaches the sensor. Approximating the collection efficiency by the fractional solid angle subtended by the imaging optics,
\begin{equation}
    \eta_{\mathrm{coll}} = \frac{r^2}{4d^2},
\end{equation}
with $r = 20\,\mathrm{mm}$ the effective lens radius and $d = 430\,\mathrm{mm}$ the distance between the cloud and the collection optics, obtaining a geometrical efficiency of:
\begin{equation}
    \eta_{\mathrm{coll}} = \frac{20^2}{4(430)^2} \approx 5.41 \times 10^{-4}.
\end{equation}

The number of detected photons per atom is
\begin{equation}
    N_{\mathrm{ph,1}} = R_{\mathrm{sc}}\,t_{\mathrm{img}}\,\eta_{\mathrm{coll}}\,T_{\mathrm{tot}},
\end{equation}
\CCn{where $T_{\mathrm{tot}}$ accounts for the transmission of the imaging system.}{Lo stesso simbolo indica due cammini ottici diversi, e la differenza e' grande. Piu' sopra $T_{\mathrm{tot}}=T_1T_{\mathrm{vac}}T_2T_3\approx0.01272$ e' la trasmissione del cammino di CALIBRAZIONE, che come dici tu stessa passa ``through a series of attenuators'': gli attenuatori servono a portare un fascio del MOT a un livello che non saturi il sensore. Qui invece serve la trasmissione del cammino di IMAGING, percorso dalla fluorescenza degli atomi, che quegli attenuatori non li ha. Se nel conto e' finito lo $0.01272$, $N_{\mathrm{ph},1}$ risulta troppo piccolo di circa il fattore degli attenuatori e il numero di atomi troppo grande dello stesso fattore, che e' dell'ordine delle decine. Se invece hai usato un valore diverso, allora e' solo il simbolo a essere riusato, ma va comunque distinto: chiamali $T_{\mathrm{cal}}$ e $T_{\mathrm{img}}$ e scrivi entrambi i valori.}

The absolute number of trapped atoms is finally obtained from
\begin{equation}
    N = \frac{N_{\mathrm{det}}}{N_{\mathrm{ph,1}}}
    = \frac{K\,C_{\mathrm{tot}}}{R_{\mathrm{sc}}\,t_{\mathrm{img}}\,\eta_{\mathrm{coll}}\,T_{\mathrm{tot}}}.
\end{equation}

\section{Expected Number of Trapped Atoms}
\label{sec:exp_n_trapped}
A simple estimate of the number of atoms loaded into the ODT can be obtained by combining the spatial density distribution of the MOT with the thermal velocity distribution of the atoms. An atom at position $\vec{r}$ is assumed to be captured if its transverse kinetic energy is smaller than the local trap depth:
\begin{equation}
    \frac{1}{2}m\left(v_x^2 + v_y^2\right) < \left|U_{\mathrm{ODT}}(\vec{r})\right|.
\end{equation}

The spatial density of the initial cloud is denoted by $n_{\mathrm{MOT}}(\vec{r})$, while the velocity distribution is taken to be Maxwell-Boltzmann,
\begin{equation}
    f(\vec{v}) = \frac{1}{(2\pi\mu^2)^{3/2}} \exp\left(-\frac{v^2}{2\mu^2}\right),
    \qquad
    \mu^2 = \frac{k_B T}{m}.
\end{equation}

The number of captured atoms is obtained by integrating over phase space, \CCn{restricting the result to atoms moving toward the trap region}{Il fattore $1/2$ mescola due criteri diversi: la CATTURA (energia cinetica trasversa sotto la profondita' locale) e la DIREZIONE del moto assiale. Un atomo catturato radialmente con $v_z>0$ e' comunque intrappolato, semplicemente va dall'altra parte; escluderlo significa che la quantita' calcolata non e' ``the number of atoms loaded into the ODT'', come dice la frase d'apertura, ma il numero di atomi caricati E diretti verso la fibra. Il capitolo 3 se ne accorge e toglie il fattore $1/2$ proprio per questo, spiegandolo in una nota. Meglio definire qui una volta sola le due quantita', con due nomi diversi, invece di correggere il fattore a valle.}, namely those with $v_z < 0$. Under this assumption,
\begin{equation}
    N_{\mathrm{ODT}} =
    \int d^3r\, n_{\mathrm{MOT}}(\vec{r})\,
    \frac{1}{2}
    \int_{v_x^2+v_y^2<\frac{2|U|}{m}}
    dv_x\,dv_y\,
    \frac{1}{2\pi\mu^2}
    \exp\left(-\frac{v_x^2+v_y^2}{2\mu^2}\right).
\end{equation}

The velocity integral is more conveniently evaluated in polar coordinates in the transverse velocity plane. Using
\begin{equation}
    dv_x\,dv_y = \pi\,d(v_r^2),
\end{equation}
and introducing the variable
\begin{equation}
    x = \frac{v_r^2}{2\mu^2},
\end{equation}
one obtains
\begin{equation}
    \int_{0}^{\frac{2|U|}{m}}
    \pi\,d(v_r^2)\,
    \frac{1}{2\pi\mu^2}
    \exp\left(-\frac{v_r^2}{2\mu^2}\right)
    =
    \int_{0}^{\frac{|U|}{m\mu^2}} e^{-x}\,dx
    =
    1 - \exp\left(-\frac{|U_{\mathrm{ODT}}(\vec{r})|}{k_B T}\right).
\end{equation}

The expression for the number of trapped atoms therefore becomes
\begin{equation}
    N_{\mathrm{ODT}} =
    \int d^3r\, n_{\mathrm{MOT}}(\vec{r})\,
    \frac{1}{2}
    \left[
        1 - \exp\left(-\frac{|U_{\mathrm{ODT}}(\vec{r})|}{k_B T}\right)
    \right].
\end{equation}

\CCn{In the shallow-trap limit, $|U_{\mathrm{ODT}}| \ll k_B T$, the exponential can be expanded to first order one then finds:}{Verificato che l'approssimazione regge, e conviene dirlo perche' a prima vista sembra il contrario. Se si guarda la profondita' al waist non torna affatto: dalla tua stessa $\eta_{\mathrm{shallow}}=0.51\%$ si ricava $U_0\approx167~\mu$K, cioe' $U_0/k_BT\approx20$, dove sviluppare $1-e^{-x}$ al primo ordine sarebbe assurdo. Ma il waist non e' dove stanno gli atomi: alla posizione del MOT il fascio e' largo $126~\mu$m, quindi la profondita' \emph{locale} scende di $(19/126)^2$ e vale circa $3.7~\mu$K, cioe' $u=U/k_BT\approx0.43$. Li' l'approssimazione e' buona: il rapporto esatto fra l'integrale vero e quello linearizzato e' $\mathrm{Ein}(u)/u=0.90$, quindi il modello sovrastima $\eta$ di circa il $10$ per cento e non di piu'. Due cose da aggiungere, entrambe di una riga: che la condizione va verificata sulla profondita' locale alla posizione di carico e non al fuoco, il che e' proprio il motivo per cui vale; e che il $10$ per cento residuo e' un sistematico noto del modello, con il suo segno, utile quando si confronta con il valore misurato. Piccola incoerenza di forma nella stessa frase: ``expanded to first order one then finds'' ha perso la punteggiatura fra le due proposizioni.}
\begin{equation}
    N_{\mathrm{ODT}}
    \simeq
    \int d^3r\, n_{\mathrm{MOT}}(\vec{r})\,
    \frac{|U_{\mathrm{ODT}}(\vec{r})|}{2k_B T}.
\end{equation}

For a Gaussian dipole trap,
\begin{equation}
    U_{\mathrm{ODT}}(\vec{r}) =
    -U_0\,
    \frac{W_0^2}{W^2(z)}
    \exp\left(-\frac{2(x^2+y^2)}{W^2(z)}\right),
\end{equation}
so that
\begin{equation}
    N_{\mathrm{ODT}}
    \simeq
    \frac{U_0}{2k_B T}
    \int d^3r\, n_{\mathrm{MOT}}(\vec{r})\,
    \frac{W_0^2}{W^2(z)}
    \exp\left(-\frac{2(x^2+y^2)}{W^2(z)}\right).
\end{equation}

If the ODT waist is much smaller than the transverse size of the MOT (if $W(z) \ll R_x, R_y$), the transverse integrals can be extended to infinity:
\begin{equation}
    \int_{-\infty}^{+\infty}\int_{-\infty}^{+\infty}
    \exp\left(-\frac{2(x^2+y^2)}{W^2(z)}\right)\,dx\,dy
    =
    \frac{\pi W^2(z)}{2}.
\end{equation}
Substituting this result into the previous expression leads to the cancellation of the $W^2(z)$ term:
\begin{equation}
    N_{\mathrm{ODT}}
    =
    \frac{U_0}{2k_B T}
    \frac{\pi W_0^2}{2}
    \int_{-R_z}^{R_z} n_{\mathrm{MOT}}(0,0,z)\,dz.
\end{equation}

To obtain a compact analytical estimate, the MOT is modeled as a uniform ellipsoid with radii $R_x$, $R_y$, and $R_z$, and constant density
\begin{equation}
    n_{\mathrm{MOT}} =
    \frac{3N_{\mathrm{MOT}}}{4\pi R_x R_y R_z}.
\end{equation}
Under this approximation,
\begin{equation}
    N_{\mathrm{ODT}}
    =
    \frac{U_0 \pi W_0^2}{4k_B T}
    \left(
        \frac{3N_{\mathrm{MOT}}}{4\pi R_x R_y R_z}\cdot 2R_z
    \right)
    =
    \frac{3}{8}
    \frac{U_0 W_0^2}{k_B T R_x R_y}
    N_{\mathrm{MOT}}.
\end{equation} 
